\documentclass[uplatex,dvipdfmx]{jlreq}
\usepackage[fleqn,tbtags]{mathtools}
\documentclass[uplatex,dvipdfmx]{jsarticle}

\usepackage[uplatex,deluxe]{otf} % UTF
\usepackage[noalphabet]{pxchfon} % must be after otf package
\usepackage{stix2} %欧文＆数式フォント
\usepackage[fleqn,tbtags]{mathtools} % 数式関連 (w/ amsmath)
\usepackage{hira-stix} % ヒラギノフォント＆STIX2 フォント代替定義（Warning回避）
\usepackage{float}
\usepackage{siunitx}

\begin{document}

\title{第二回 事前課題: 目標を達成するための案出し}
\author{25G1059 佐藤涼菜}
\date{2026年4月22日}
\maketitle




\section{はじめに}

\subsection{社会的背景}
人間は音楽を聞く際に，聴取者の音楽経験や学習に基づいて次の音を予期する音楽的期待という感情反応を示す．
特に万人が一度は聞いたことのあるような輪唱は次の音を予期することが容易である．この音楽的期待は予期された展開通りであれば聴取者に喜びや満足感を与えるが，
予期されない展開では驚きや不安など負の感情を喚起する\cite{ref:kitai}．したがってArduino間の同期にズレが生じるほど，聴取者は予期していた展開ではない演奏
により負の感情を想起させてしまう問題が生じる．

\subsection{技術的背景}
聴取者の音楽的期待を裏切らないように楽器間の同期のズレをなくすには，許容される通信誤差の範囲を設定する必要がある．
Arduinoオーケストラでズレたように聞こえる，すなわち予期しない離散的反射音(エコー)が生じる遅延時間は50 msを超える場合
であることが先行研究から示されている\cite{ref:eko}．よって，制作する楽器間同期演奏システムの許容される通信遅延は50 ms以内を
最低目標とする．




\section{初回のハッカソンで決まった班の方針}

\subsection{構成}
・楽器3個，リズム1個，楽曲表現用1個程度，通信用1個\par
・楽器について:\par
リズムはドラムで表現する．また，メロディーに関して，現時点でピアノは確定でありその他楽器を２個選択する予定である．
第二回のハッカソンでは，残りの楽器の案出しを行う．

\subsection{表現方法例}
前回出た表現方法は以下のとおりである．\par
・LEDを用いて音の周波数ごとに光るLEDを設定する．\par
・サーボモーターを使って腕が動く．\par
・bpmの変更で観覧車の回転速度を変える．\par
・音程をMatrix上に流す．\par

\subsection{通信方法例}
まず，楽器間の通信方法はwifiによる無線通信が第一候補に挙がった．
また楽器の同期の方法としては，入力で加速度センサを用いてロボットが同期ボタンを押すことで演奏がされるようにする案が出た．
第二回のハッカソンでは，主にこの同期の通信について具体的な方法の案出しを行う．


\section{既存のシステム}

% \subsection{構成}
% まず，演奏を構成する楽器について案出しを行うためにprocessingで楽器の音色を表現する技術について述べる．

% \subsection{通信方法}
% 無線通信の基本的なプロトコルとして，TCPとUDPが存在する．TCPとはTransmission Control Protocolの略でパケットの欠落や重複を防ぎ信頼性の高い通信を実現することが可能である．
% したがって，今回のテーマでは正確な楽器音を流すことが期待できる．ただし，パケットが失われた場合に再送する手間が生じてしまうため，リアルタイム性は低い．
% 対して，UDPはでデータを送信する際に接続の確立や到達確認を行わなずリアルタイム性が高いことが特徴である．よって今回のテーマでは演奏の同期するタイミングの
% 誤差が短くすることが期待できる．

% 既存のネットワーク型音楽演奏システムとして，YAMAHAの「SYNCROOM」が挙げられる\cite{ref:yamaha}．このシステムは，距離が離れていてもリアルタイムに音楽演奏が可能であるアプリケーションであり，
% YAMAHAが独自に開発した，インターネット回線を介してオーディオデータの双方向送受信を遅れを最小化するための「NETDUETTO」を用いて演奏合奏を可能にしている．



\section{自身のアイディア}

\subsection{構成}
表現する楽器に関しては，より楽器音らしい演奏を行うために，倍音成分が少なく音響合成の基本である正弦波に近い楽器，すなわちリコーダーやフルートを提案する．
倍音成分が多く，複雑な合成が必要な金管楽器を選択してしまうと，CPUへの負荷が高まり演奏の同期に誤差が生じてしまう可能性がある．そのため，同期のタイミングの誤差を最小化し受信精度を向上させ
るためにシンプルな正弦波で表現しやすい楽器を選択した．
また，もう一個の楽器としてADSR制御が単純化しやすい木琴を提案する．
ADSR制御とは，音量や音色などの時間的な変化を制御することである\cite{ref:adsr}．
ADSRとは，それぞれがコントロールするパラメータの頭文字であり以下に示す．

・ATTACK TIME:音が発生してから最大レベルに達するまでの時間．\par
・DECAY TIME：アタック後サスティンレベルまでに達するまでの時間．\par
・SUSTAIN LEVEL：音が維持されている際の音量のレベル．\par
・RELEASE TIME：音を止める操作をした後，完全に消えるまでの時間．\par

ADSR制御が単純であれば演奏タイミングの誤差の最小化につながり，また楽器音らしさの再現度は高くなることが予想される．
そこで，私が提案するのは木琴である．木琴は音盤を叩いてから最大の音になるまでの時間が短く，また叩いた後はすぐに音が消えるため
比較的ADSR制御の処理が少ないと考えられる．

したがって，楽器音らしい音をズレることなく同期させることを第一の目標とすると，残りの2個の楽器は
フルートやリコーダーと木琴が適していると考える．

\subsection{通信方法}
次に，楽器間の同期を行う通信方法に関しては，UDPを用いた無線通信を提案する．無線通信の基本的なプロトコルとして，TCPとUDPが存在する．TCPとはパケットの欠落や重複を防ぎ信頼性の高い通信を実現することが可能である．
ただし，パケットが失われた場合に再送する手間が生じてしまうため，リアルタイム性は低い．対して，UDPはデータを送信する際に接続の確立や到達確認を行わずリアルタイム性が高いことが特徴である．よって今回のテーマでは演奏の同期するタイミングの
誤差が短くすることが期待できる．また，演奏の開始やテンポが変化した際に楽器間でズレが生じないようフィルターを設け余計な音を除去する仕様を提案する．

















\begin{thebibliography}{999}

\bibitem{ref:kitai}星野悦子．音楽心理学の研究動向ー音楽で喚起される感情の研究を中心にー．音楽教育学第51-2.(2022)．
\bibitem{ref:eko}Helmut Haas．Über den Einfluss eines Einfachechos auf die Hörsamkeit von Sprache．Acustica, Vol. 1, pp. 49-58．
\bibitem{ref:yamaha}YAMAHA.SYNCROOM公式HP．
\small\url{https://syncroom.yamaha.com/}

\bibitem{ref:adsr}島村．アナログシンセ超入門～その2：EG（エンベロープジェネレーター、ADSR）編.(2026年4月21日閲覧)
\small\url{https://info.shimamura.co.jp/digital/guide/2018/02/122103}

\end{thebibliography}

\end{document}




